\documentclass[11pt,a4paper]{article}
\usepackage{fontspec}
\setmainfont{DejaVu Sans}
\setsansfont{DejaVu Sans}
\usepackage{polyglossia}
\setdefaultlanguage{russian}
\setotherlanguage{english}
\usepackage{amsmath,amssymb,mathtools}
\usepackage{geometry}
\usepackage{xcolor}
\usepackage{enumitem}
\usepackage{microtype}
\usepackage{fancyhdr}
\usepackage{tcolorbox}
\usepackage{hyperref}
\geometry{left=18mm,right=18mm,top=18mm,bottom=20mm}
\setlength{\parindent}{0pt}
\setlength{\parskip}{4pt}
\definecolor{hseblue}{HTML}{003D7C}
\definecolor{softblue}{HTML}{EEF4FA}
\hypersetup{colorlinks=true,linkcolor=hseblue,urlcolor=hseblue}
\pagestyle{fancy}
\fancyhf{}
\lhead{Методы оптимизации в машинном обучении}
\rhead{Семинар 4}
\cfoot{\thepage}
\newcommand{\R}{\mathbb{R}}
\newcommand{\norm}[1]{\left\lVert #1\right\rVert}
\newcommand{\grad}{\nabla}
\newcounter{problem}
\newenvironment{problem}[1][]{%
  \refstepcounter{problem}%
  \begin{tcolorbox}[colback=white,colframe=hseblue!70!black,boxrule=0.7pt,arc=2mm,left=2.5mm,right=2.5mm,top=2mm,bottom=2mm,title={\textbf{Задача \theproblem.} #1},fonttitle=\normalsize]
}{\end{tcolorbox}}
\begin{document}

{\LARGE\bfseries\color{hseblue} Семинар 4. Гладкость, выпуклость и условия оптимальности}

\vspace{0.4em}
Материал: вторая половина лекции 4 и лекция 5. Основные задачи рекомендуется разобрать на семинаре; дополнительные --- использовать как резерв и для самостоятельной работы.

\section*{Основные задачи}

\begin{problem}[$L$-гладкость и гессиан]
Рассмотрим функцию
\[
f(x,y)=\frac12x^2+2y^2+\sin x.
\]
\begin{enumerate}[label=\alph*)]
\item Найдите $\grad f(x,y)$ и $\grad^2 f(x,y)$.
\item Найдите наименьшую константу $L$, для которой
\[
\|\grad^2 f(x,y)\|_2\le L
\qquad\text{для всех }(x,y)\in\R^2.
\]
\item Используя найденную константу, объясните, почему
\[
\|\grad f(u)-\grad f(v)\|_2\le L\|u-v\|_2
\qquad\text{для любых }u,v\in\R^2.
\]
\item Является ли функция выпуклой? Сильно выпуклой на всей $\R^2$?
\end{enumerate}
\end{problem}

\begin{problem}[Лемма об убывании и шаг градиентного спуска]
Пусть $f$ --- дифференцируемая $L$-гладкая функция и известно, что
\[
f(x+h)\le f(x)+\grad f(x)^\top h+\frac L2\|h\|_2^2.
\]
\begin{enumerate}[label=\alph*)]
\item Подставьте $h=-\alpha\grad f(x)$ и выведите оценку
\[
f(x-\alpha\grad f(x))
\le
f(x)-\alpha\left(1-\frac{\alpha L}{2}\right)\|\grad f(x)\|_2^2.
\]
\item При каких $\alpha>0$ эта оценка гарантирует строгое уменьшение значения функции, если $\grad f(x)\neq0$?
\item Что получается при $\alpha=1/L$?
\item Верно ли утверждение: «если $\alpha>2/L$, то значение функции обязательно возрастёт»? Обоснуйте ответ.
\end{enumerate}
\end{problem}

\begin{problem}[Выпуклость по определению и неравенство Йенсена]
Для каждой функции определите, выпукла ли она на указанном множестве:
\[
f_1(x)=x^2,\quad x\in\R;
\qquad
f_2(x)=|x|,\quad x\in\R;
\]
\[
f_3(x)=x^3,\quad x\in\R;
\qquad
f_4(x)=-\log x,\quad x>0.
\]
\begin{enumerate}[label=\alph*)]
\item Для $f_1$ докажите выпуклость непосредственно из определения, не используя производные. Полезно рассмотреть разность
\[
(1-t)x^2+ty^2-\bigl((1-t)x+ty\bigr)^2.
\]
\item Для $f_2,f_3,f_4$ дайте доказательство или контрпример.
\item Используя выпуклость $-\log x$ и неравенство Йенсена с весами $1/2,1/2$, получите неравенство между средним арифметическим и средним геометрическим двух положительных чисел.
\end{enumerate}
\end{problem}

\begin{problem}[Критерий первого порядка и глобальный минимум]
Рассмотрим функцию
\[
f(x,y)=e^x+\frac12y^2.
\]
\begin{enumerate}[label=\alph*)]
\item Найдите $\grad f(x,y)$.
\item Для точки $x_0=(0,1)$ выпишите линейную модель
\[
m_{x_0}(z)=f(x_0)+\grad f(x_0)^\top(z-x_0).
\]
\item Докажите, что $f(z)\ge m_{x_0}(z)$ для любого $z\in\R^2$.
\end{enumerate}

Теперь рассмотрим
\[
g(x,y)=x^2+2y^2-2x+4y.
\]
\begin{enumerate}[label=\alph*),resume]
\item Найдите стационарную точку $g$.
\item Не сравнивая её со всеми остальными точками, докажите, что это глобальный минимум.
\end{enumerate}
\end{problem}

\begin{problem}[Гессиан, сильная выпуклость и параметры $\mu,L$]
Рассмотрим
\[
f(x,y)=x^2+xy+2y^2+\frac14x^4.
\]
\begin{enumerate}[label=\alph*)]
\item Найдите гессиан $\grad^2f(x,y)$.
\item Докажите, что функция выпукла на всей $\R^2$.
\item Найдите максимально возможную константу $\mu>0$, для которой
\[
\grad^2f(x,y)\succeq \mu I
\qquad\text{для всех }(x,y)\in\R^2.
\]
\item Является ли функция глобально $L$-гладкой на всей $\R^2$?
\item Найдите подходящую константу $L$ на полосе $|x|\le1$.
\end{enumerate}
\end{problem}

\section*{Дополнительные задачи}

\begin{problem}[Почему локальный минимум выпуклой функции глобален]
Пусть $f$ выпукла на выпуклом множестве $X$, а $x^\star\in X$ --- локальный минимум. Докажите, что $x^\star$ является глобальным минимумом.

\textit{Указание.} Предположите обратное: существует $y\in X$ такое, что $f(y)<f(x^\star)$. Рассмотрите
\[
x_t=(1-t)x^\star+ty
\]
при достаточно малом $t>0$.
\end{problem}

\begin{problem}[Строгая и сильная выпуклость]
Для функций
\[
f_1(x)=x^2,
\qquad
f_2(x)=x^4,
\qquad
f_3(x)=e^x
\]
на $\R$ определите:
\begin{enumerate}[label=\alph*)]
\item какие из них выпуклы;
\item какие строго выпуклы;
\item какие сильно выпуклы на всей $\R$;
\item какие становятся сильно выпуклыми на отрезке $[-R,R]$, $R>0$.
\end{enumerate}
Также докажите, что строго выпуклая функция имеет не более одного глобального минимума.
\end{problem}

\begin{problem}[Операции, сохраняющие выпуклость]
Пусть $f_1,\ldots,f_m$ --- выпуклые функции и $\alpha_i\ge0$.
\begin{enumerate}[label=\alph*)]
\item Докажите, что
\[
F(x)=\sum_{i=1}^m\alpha_i f_i(x)
\]
выпукла.
\item Определите, выпукла ли функция
\[
F_1(x)=e^x+x^2+|x|.
\]
\item Можно ли из пункта a) сделать вывод о выпуклости
\[
F_2(x)=x^2-3e^x?
\]
Если нет, исследуйте её выпуклость отдельно.
\item Приведите пример, показывающий, что наличие отрицательного коэффициента в линейной комбинации выпуклых функций не обязательно уничтожает выпуклость.
\end{enumerate}
\end{problem}

\begin{problem}[Ridge как источник сильной выпуклости]
Рассмотрим
\[
f(w)=\frac12\|Xw-y\|_2^2+\frac\lambda2\|w\|_2^2,
\qquad \lambda>0.
\]
\begin{enumerate}[label=\alph*)]
\item Найдите гессиан и докажите
\[
\grad^2f(w)=X^\top X+\lambda I\succeq\lambda I.
\]
\item Объясните, почему задача становится сильно выпуклой даже при вырожденной матрице $X^\top X$.
\item Пусть собственные значения $X^\top X$ равны $0,1,20$. Найдите собственные значения гессиана Ridge-задачи и её число обусловленности как функцию $\lambda$.
\item Исследуйте пределы полученного числа обусловленности при $\lambda\to0+$ и $\lambda\to\infty$.
\end{enumerate}
\end{problem}

\begin{problem}[Выпуклость логистической регрессии]
Рассмотрим логистическую потерю
\[
\ell(z)=\log(1+e^{-z}).
\]
\begin{enumerate}[label=\alph*)]
\item Найдите $\ell'(z)$ и $\ell''(z)$.
\item Докажите, что $\ell$ выпукла.
\item Объясните, почему для фиксированных $x\in\R^d$ и $y\in\{-1,1\}$ функция
\[
w\mapsto \ell(yx^\top w)
\]
выпукла по $w$.
\item Почему сумма таких потерь по объектам остаётся выпуклой?
\item Обязана ли задача логистической регрессии без регуляризации быть строго выпуклой? Сильно выпуклой на всей $\R^d$? Какие свойства данных здесь существенны?
\item Что меняется после добавления регуляризатора $\frac\lambda2\|w\|_2^2$, $\lambda>0$?
\end{enumerate}
\end{problem}

\end{document}